\chapter*{Executive summary}

\noindent
\begin{longtable}{|p{\dimexpr \linewidth-2\tabcolsep-2\arrayrulewidth}|}
	
\hline%------------------------------------------------------------
\sumheading  Title of Project \\
\hline%------------------------------------------------------------
 Joint Design for a Portable Coordinate Measurement Arm (PCMA) \\[1ex]

\hline%------------------------------------------------------------
\sumheading  Objectives \\
\hline%------------------------------------------------------------
 To derive technical performance measures (TPMs) from a generalised error-propagation model. To design, test, and build a joint with position-holding implementation that will meet the derived TPMs and maintain and/or improve the arm's accuracy.\\[1ex]

\hline%------------------------------------------------------------
\sumheading  What is current practice and what are its limitations? \\
\hline%------------------------------------------------------------
 Current joint use counterbalancing and have little to no holding torque. The PCMAs are reliant on the operator to support the arm during operation. \\[1ex]

\hline%------------------------------------------------------------
\sumheading  What is new in this project? \\
\hline%------------------------------------------------------------
 The proposed joint design will implement position-holding mechanism designed from the TPMs that were derived from the error-propagation model.\\[1ex]

\hline%------------------------------------------------------------
\sumheading  If the project is successful, how will it make a difference? \\
\hline%------------------------------------------------------------
 Implementation of the joint design will maintain and/or improve the measurement accuracy of the arm and a modified kinematic model of the PCMA will be derived to create a more generalised framework for accuracy analysis.\\[1ex]

\hline%------------------------------------------------------------
\sumheading  What are the risks to the project being a success? Why is it expected to be successful? \\
\hline%------------------------------------------------------------
 Manufacturing tolerances and the M\&M workshop's capabilities might not meet derived requirements. The position-holding mechanism might introduce unexpected friction that impedes the free movement of the PCMA. The project is expected to succeed because the requirements are based off of an established error model.\\[1ex]

\hline%------------------------------------------------------------
\sumheading  What contributions have/will other students made/make? \\
\hline%------------------------------------------------------------
 ... \\[1ex]

\hline%------------------------------------------------------------
\sumheading  Which aspects of the project will carry on after completion and why? \\
\hline%------------------------------------------------------------
 ... \\[1ex]

\hline%------------------------------------------------------------
\sumheading  What arrangements have been/will be made to expedite continuation? \\
\hline%------------------------------------------------------------
 ... \\[1ex]

\hline%------------------------------------------------------------
\end{longtable}

